\setcounter{ex}{0}

% --- CÂU 1 ---
\begin{ex}
	Cho $\cos 2\alpha = \dfrac{1}{8}$. Giá trị của $\sin^2\alpha$ bằng
	\choice
	{$\dfrac{1}{16}$}
	{\True $\dfrac{7}{16}$}
	{$\dfrac{9}{16}$}
	{$\dfrac{7}{8}$}
	\loigiai{
		Áp dụng công thức hạ bậc $\sin^2\alpha = \dfrac{1 - \cos 2\alpha}{2}$, ta có:
		$$\sin^2\alpha = \dfrac{1 - \frac{1}{8}}{2} = \dfrac{7}{16}.$$
	}
\end{ex}

% --- CÂU 2 ---
\begin{ex}
	Cho $\tan a = \dfrac{1}{2}$ và $\tan b = \dfrac{1}{3}$. Giá trị của $\tan(a - b)$ bằng
	\choice
	{\True $\dfrac{1}{7}$}
	{$\dfrac{1}{5}$}
	{$7$}
	{$\dfrac{5}{7}$}
	\loigiai{
		Ta có:
		$$\tan(a - b) = \dfrac{\tan a - \tan b}{1 + \tan a\tan b} = \dfrac{\frac{1}{2} - \frac{1}{3}}{1 + \frac{1}{2} \cdot \frac{1}{3}} = \dfrac{\frac{1}{6}}{\frac{7}{6}} = \dfrac{1}{7}.$$
	}
\end{ex}

% --- CÂU 3 ---
\begin{ex}
	Tập xác định của hàm số $y = \dfrac{2}{\sin x + 1}$ là
	\choice
	{$\mathbb{R} \setminus \{k\pi,\ k \in \mathbb{Z}\}$}
	{$\mathbb{R} \setminus \left\{-\dfrac{\pi}{2} + k\pi,\ k \in \mathbb{Z}\right\}$}
	{\True $\mathbb{R} \setminus \left\{-\dfrac{\pi}{2} + k2\pi,\ k \in \mathbb{Z}\right\}$}
	{$\mathbb{R} \setminus \left\{\dfrac{\pi}{2} + k2\pi,\ k \in \mathbb{Z}\right\}$}
	\loigiai{
		Hàm số xác định khi và chỉ khi:
		$$\sin x + 1 \ne 0 \iff \sin x \ne -1 \iff x \ne -\dfrac{\pi}{2} + k2\pi\ (k \in \mathbb{Z}).$$
		Vậy tập xác định của hàm số là $\mathscr{D} = \mathbb{R} \setminus \left\{-\dfrac{\pi}{2} + k2\pi,\ k \in \mathbb{Z}\right\}$.
	}
\end{ex}

% --- CÂU 4 ---
\begin{ex}
	Chu kì tuần hoàn của hàm số $y = \tan\left(2x + \dfrac{\pi}{3}\right)$ là
	\choice
	{\True $\dfrac{\pi}{2}$}
	{$2\pi$}
	{$\dfrac{\pi}{4}$}
	{$\pi$}
	\loigiai{
		Hàm số $y = \tan\left(2x + \dfrac{\pi}{3}\right)$ tuần hoàn với chu kì $T = \dfrac{\pi}{2}$.
	}
\end{ex}

% --- CÂU 5 ---
\begin{ex}
	Nghiệm của phương trình $\cot\left(x + \dfrac{\pi}{6}\right) = -\sqrt{3}$ là
	\choice
	{$x = \dfrac{\pi}{6} + k\pi\ (k \in \mathbb{Z})$}
	{$x = \dfrac{\pi}{3} + k\pi\ (k \in \mathbb{Z})$}
	{$x = -\dfrac{\pi}{6} + k\pi\ (k \in \mathbb{Z})$}
	{\True $x = -\dfrac{\pi}{3} + k\pi\ (k \in \mathbb{Z})$}
	\loigiai{
		Ta có:
		$$\cot\left(x + \dfrac{\pi}{6}\right) = -\sqrt{3} \iff \cot\left(x + \dfrac{\pi}{6}\right) = \cot\left(-\dfrac{\pi}{6}\right) \iff x + \dfrac{\pi}{6} = -\dfrac{\pi}{6} + k\pi \iff x = -\dfrac{\pi}{3} + k\pi\ (k \in \mathbb{Z}).$$
	}
\end{ex}

% --- CÂU 6 ---
\begin{ex}
	Số nghiệm của phương trình $2\sin 2x = 1$ trên khoảng $(0; \pi)$ là
	\choice
	{$4$}
	{$3$}
	{$1$}
	{\True $2$}
	\loigiai{
		Ta có:
		$$2\sin 2x = 1 \iff \sin 2x = \dfrac{1}{2} \iff \hoac{2x &= \dfrac{\pi}{6} + k2\pi \\ 2x &= \dfrac{5\pi}{6} + k2\pi} \iff \hoac{x &= \dfrac{\pi}{12} + k\pi \\ x &= \dfrac{5\pi}{12} + k\pi}\ (k \in \mathbb{Z}).$$
		Trên khoảng $(0; \pi)$, phương trình có $2$ nghiệm là $x = \dfrac{\pi}{12}$ và $x = \dfrac{5\pi}{12}$ (ứng với $k = 0$).
	}
\end{ex}

% --- CÂU 7 ---
\begin{ex}
	Giá trị lớn nhất của hàm số $y = \sqrt{3}\sin x - \cos x$ bằng
	\choice
	{$\sqrt{3}$}
	{$4$}
	{\True $2$}
	{$1 + \sqrt{3}$}
	\loigiai{
		Ta có $y = \sqrt{3}\sin x - \cos x = 2\sin\left(x - \dfrac{\pi}{6}\right)$.\\
		Vì $-1 \le \sin\left(x - \dfrac{\pi}{6}\right) \le 1, \forall x \in \mathbb{R}$ nên $-2 \le y \le 2$.\\
		Vậy giá trị lớn nhất của hàm số bằng $2$.
	}
\end{ex}

% --- CÂU 8 ---
\begin{ex}
	Cho cấp số cộng $(u_n)$ thỏa mãn $u_1 + u_2 + u_3 = -9$ và $u_5 = 6$. Số hạng tổng quát của cấp số cộng là
	\choice
	{\True $u_n = 3n - 9$}
	{$u_n = 3n - 6$}
	{$u_n = -3n + 3$}
	{$u_n = 2n - 6$}
	\loigiai{
		Gọi $d$ là công sai của cấp số cộng. Ta có hệ phương trình:
		$$\heva{u_1 + u_2 + u_3 &= -9 \\ u_5 &= 6} \iff \heva{3u_1 + 3d &= -9 \\ u_1 + 4d &= 6} \iff \heva{u_1 &= -6 \\ d &= 3.}$$
		Số hạng tổng quát của cấp số cộng là $u_n = u_1 + (n - 1)d = -6 + 3(n - 1) = 3n - 9$.
	}
\end{ex}

% --- CÂU 9 ---
\begin{ex}
	Cho cấp số cộng $(u_n)$ có $u_1 = 3$ và công sai $d = -2$. Số hạng đầu tiên nhỏ hơn $-50$ là số hạng thứ
	\choice
	{$27$}
	{\True $28$}
	{$26$}
	{$29$}
	\loigiai{
		Ta có $u_n = u_1 + (n - 1)d = 3 - 2(n - 1) = 5 - 2n$.\\
		Xét $u_n < -50 \iff 5 - 2n < -50 \iff 2n > 55 \iff n > 27{,}5$.\\
		Vì $n \in \mathbb{N}^*$ nên số hạng đầu tiên nhỏ hơn $-50$ là số hạng thứ $28$.
	}
\end{ex}

% --- CÂU 10 ---
\begin{ex}
	Cho hình chóp $S.ABCD$ có đáy $ABCD$ là hình thang ($AB \parallel CD,\ AB = 2CD$), $O = AC \cap BD$, $M, N$ lần lượt là trung điểm của $SA, SB$ (hình vẽ dùng chung cho câu 10, 11, 12). Giao tuyến của hai mặt phẳng $(SAB)$ và $(SCD)$ là đường thẳng đi qua $S$ và
	\begin{center}
		\begin{tikzpicture}[scale=0.85, font=\footnotesize, line join=round, line cap=round]
			\coordinate (A) at (0,0);
			\coordinate (B) at (4,0);
			\coordinate (D) at (1,1.1);
			\coordinate (C) at (3,1.1);
			\coordinate (O) at (intersection of A--C and B--D);
			\coordinate (S) at (2,3.5);
			\coordinate (M) at ($(S)!0.5!(A)$);
			\coordinate (N) at ($(S)!0.5!(B)$);
			\draw[thick] (S)--(A) (S)--(B) (A)--(B) (M)--(N);
			\draw[dashed] (S)--(D) (S)--(C) (A)--(D)--(C)--(B) (A)--(C) (B)--(D);
			\foreach \p/\pos in {S/above, A/below left, B/below right, C/right, D/left, O/below, M/above left, N/above right}
				\node[\pos] at (\p) {$\p$};
			\foreach \p in {S,A,B,C,D,O,M,N}
				\fill (\p) circle (1.2pt);
		\end{tikzpicture}
	\end{center}
	\choice
	{song song với $BC$}
	{song song với $AD$}
	{đi qua $O$}
	{\True song song với $AB$}
	\loigiai{
		Hai mặt phẳng $(SAB)$ và $(SCD)$ có điểm chung là $S$ và lần lượt chứa hai đường thẳng song song $AB, CD$ ($AB \parallel CD$).\\
		Do đó giao tuyến của hai mặt phẳng $(SAB)$ và $(SCD)$ là đường thẳng đi qua $S$ và song song với $AB$ (và $CD$).
	}
\end{ex}

% --- CÂU 11 ---
\begin{ex}
	Gọi $F$ là giao điểm của đường thẳng $SC$ và mặt phẳng $(ADN)$. Tỉ số $\dfrac{SF}{SC}$ bằng
	\choice
	{\True $\dfrac{2}{3}$}
	{$\dfrac{1}{3}$}
	{$\dfrac{3}{4}$}
	{$\dfrac{1}{2}$}
	\loigiai{
		Trong mặt phẳng $(ABCD)$, gọi $E = AD \cap BC$. Khi đó $(ADN) \cap (SBC) = EN$, suy ra $F = SC \cap EN$.\\
		Trong tam giác $ABE$, vì $CD \parallel AB$ và $CD = \dfrac{1}{2}AB$ nên $C$ là trung điểm của $BE$.\\
		Xét tam giác $SBE$ có hai đường trung tuyến $SC$ và $EN$ cắt nhau tại $F$, nên $F$ là trọng tâm của tam giác $SBE \implies \dfrac{SF}{SC} = \dfrac{2}{3}$.
	}
\end{ex}

% --- CÂU 12 ---
\begin{ex}
	Khẳng định nào sau đây đúng?
	\choice
	{$MN \parallel (SAB)$}
	{\True $MN \parallel (SCD)$}
	{$MN \parallel (SBC)$}
	{$MN \parallel (SAD)$}
	\loigiai{
		Vì $M, N$ lần lượt là trung điểm của $SA, SB$ nên $MN$ là đường trung bình của tam giác $SAB \implies MN \parallel AB$.\\
		Mà $AB \parallel CD$ nên $MN \parallel CD$. Do $CD \subset (SCD)$ và $MN \not\subset (SCD)$ nên $MN \parallel (SCD)$.
	}
\end{ex}
